% SEGMENT OLP-0045-S01
% Part: sets-functions-relations
% Chapter: arithmetization
% Section: cuts
%
\documentclass[../../../include/open-logic-section]{subfiles}

\begin{document}

\olfileid{sfr}{arith}{cuts}
\olsection{$\Rat$ இலிருந்து $\Real$ க்கு}

சாரத்தில், மெய்யெண் கோட்டின் எந்தப் புள்ளி $\alpha$ உம்
அந்தக் கோட்டை மிகச்சரியாக இரு பாதிகளாகப் பிரிக்கிறது என்பதை
முழுமைப் பண்பு காட்டுகிறது: ஒன்றிற்கு $\alpha$ மீச்சிறு மேல்வரம்பாகவும்,
மற்றொன்றிற்கு $\alpha$ மீப்பெரு கீழ்வரம்பாகவும் உள்ளது.
விகிதமுறு எண்களிலிருந்து மெய்யெண்களைக் \emph{கட்டமைக்க},
விகிதமுறு எண்களைப் பிரிக்கும் \emph{வெட்டுகளாகவே} மெய்யெண்களைக்
கருதலாம் என்று டெடெகிண்ட் முன்மொழிந்தார்.
அதாவது, $< \sqrt{2}$ என அமையும் விகிதமுறு எண்களை
$> \sqrt{2}$ என அமையும் விகிதமுறு எண்களிலிருந்து பிரிக்கும்
\emph{வெட்டுடன்} $\sqrt{2}$ ஐ அடையாளப்படுத்துகிறோம்.

இதனை ஒழுங்குபடுத்துவோம்.
விகிதமுறு எண்களை இரு பாதிகளாக வெட்டினால்,
நாம் உருவாக்கிய பிரிப்பை அதன் \emph{கீழ்ப்} பாதியை மட்டும் கொண்டு
தனித்துவமாக அடையாளப்படுத்தலாம்.
எனவே இதனைத் துல்லியமாகப் பின்வருமாறு வரையறுக்கிறோம்:

% SEGMENT OLP-0045-S02
\begin{defn}[வெட்டு]
ஒரு \emph{வெட்டு} $\alpha$ என்பது,
மீப்பெரு உறுப்பு இல்லாத, விகிதமுறு எண்களின் எந்தவொரு
வெற்றற்ற தகு தொடக்கத் துண்டும் ஆகும்.
அதாவது, $\alpha$ ஒரு வெட்டாக இருக்கும்போது,
அப்போதும் அப்போது மட்டுமே:
\begin{enumerate}
	\item \emph{வெற்றற்றது, தகுவானது}: $\emptyset \neq \alpha \subsetneq \Rat$
	\item \emph{தொடக்கத் துண்டு}: அனைத்து $p,q \in \Rat$ க்கும்,
    $p < q \in \alpha$ எனில் $p \in \alpha$
	\item \emph{மீப்பெரு உறுப்பு இல்லை}: அனைத்து $p \in \alpha$ க்கும்,
    $p < q$ என அமையும் ஒரு $q \in \alpha$ உள்ளது
\end{enumerate}
பின்னர் $\Real$ என்பது வெட்டுகளின் கணமாகும்.
\end{defn}

% SEGMENT OLP-0045-S03
இப்போது $\sqrt{2} = \Setabs{p \in \Rat}{p^2 < 2\text{
அல்லது }p < 0}$ என்று கூறலாம்.
நிச்சயமாக, இது \emph{ஒரு வெட்டுதான்} எனச் சரிபார்க்க வேண்டும்;
அதனை \olref[check]{sec} இல் தருகிறோம்.

% SEGMENT OLP-0045-S04
முன்புபோல, சில பொருள்களை வரையறுத்த பிறகு,
அவற்றின் மீதான அடிப்படைச் சார்புகளையும் தொடர்புகளையும்
அடுத்து வரையறுக்க வேண்டும். எளிய ஒன்றிலிருந்து தொடங்குகிறோம்:
\begin{align*}
	\alpha \leq \beta \text{ ஆக இருக்கும்போது, அப்போதும் அப்போது மட்டுமே, }\alpha \subseteq \beta
\end{align*}
வரிசையின் இந்த வரையறை, வெட்டுகளின் கணம் முழுமைப் பண்பைக் கொண்டுள்ளது
என்ற மைய முடிவைக் \emph{கூற} அனுமதிக்கிறது.
முழுமையாக விரித்தால், கூற்றின் வடிவம் இதுவாகும்.
$S$ என்பது மேல்வரம்பைக் கொண்ட வெற்றற்ற வெட்டுகளின் கணம் எனில்,
$S$ ஒரு மீச்சிறு மேல்வரம்பைக் கொண்டுள்ளது.
மேலும் விரிவாக: $S$ இன் மேல்வரம்பாக உள்ள ஒரு வெட்டு $\lambda$ உள்ளது;
அதாவது $(\forall \alpha \in S)\alpha \subseteq
\lambda$; மேலும் $\lambda$ இத்தகைய மீச்சிறு வெட்டு;
அதாவது $(\forall \beta \in \Real)((\forall \alpha \in S)\alpha
\subseteq \beta \lif \lambda \subseteq \beta)$.
இந்த முடிவின் நிறுவல் இதோ:

% SEGMENT OLP-0045-S05
\begin{thm}\ollabel{realcompleteness}
வெட்டுகளின் கணம் முழுமைப் பண்பைக் கொண்டுள்ளது.
\end{thm}

% SEGMENT OLP-0045-S06
\begin{proof}
$S$ என்பது மேல்வரம்பைக் கொண்ட எந்தவொரு வெற்றற்ற வெட்டுகளின்
கணமாகவும் இருக்கட்டும். $\lambda = \bigcup S$ என்க.
%\Setabs{p \in \Rat}{(\exists \alpha \in S)p \in \alpha}$$
முதலில் $\lambda$ ஒரு வெட்டு என வாதிடுகிறோம்:
\begin{enumerate}
\item $S$ வெற்றற்றதால், குறைந்தது ஒரு வெட்டு $S$ இல் உள்ளது;
எனவே $\emptyset \neq \lambda$.
$S$ என்பது வெட்டுகளின் கணம் என்பதால், $\lambda \subseteq \Rat$.
$S$ க்கு ஒரு மேல்வரம்பு இருப்பதால்,
ஒவ்வொரு வெட்டு $\alpha \in S$ இலிருந்தும் விடுபட்ட ஏதோ ஒரு
$p \in \Rat$ உள்ளது.
எனவே $p\notin \lambda$; ஆதலால் $\lambda \subsetneq \Rat$.
\item $p < q \in \lambda$ எனக் கொள்க.
அப்போது $q \in \alpha$ என அமையும் ஏதோ ஒரு $\alpha \in S$ உள்ளது.
$\alpha$ ஒரு வெட்டு என்பதால், $p \in \alpha$.
எனவே $p \in \lambda$.
\item $p \in \lambda$ எனக் கொள்க.
அப்போது $p \in \alpha$ என அமையும் ஏதோ ஒரு $\alpha \in S$ உள்ளது.
$\alpha$ ஒரு வெட்டு என்பதால், $p < q$ என அமையும் ஏதோ ஒரு
$q \in \alpha$ உள்ளது. எனவே $q \in \lambda$.
\end{enumerate}
இது வாதத்தை நிறுவுகிறது.
மேலும், $(\forall \alpha \in S)\alpha
\subseteq \bigcup S = \lambda$ என்பது தெளிவு;
அதாவது $\lambda$ என்பது $S$ இன் ஒரு மேல்வரம்பு.
இப்போது $\beta \in \mathbb{R}$ உம் ஒரு மேல்வரம்பு எனக் கொள்க;
அதாவது $(\forall \alpha \in S)\alpha \subseteq \beta$.
எந்த $p \in \Rat$ க்கும், $p \in \lambda$ எனில்,
$p \in \alpha$ என அமையும் ஒரு $\alpha \in S$ உள்ளது;
எனவே $p \in \beta$.
பொதுமைப்படுத்தினால், $\lambda \subseteq \beta$.
ஆகவே $\lambda$ என்பது $S$ இன் \emph{மீச்சிறு} மேல்வரம்பு.
\end{proof}

% SEGMENT OLP-0045-S07
எனவே முழுமைப் பண்பை நிறைவேற்றும் பல பொருள்கள் நமக்குக் கிடைத்துள்ளன.
இதனை மற்றொரு வகையில் சொன்னால், நமது வெட்டுகளில் ``இடைவெளிகள்'' இல்லை.
(எனவே, விகிதமுறு எண்களுக்குப் பதிலாக மெய்யெண்களில் மேலும்
``வெட்டுகளை'' எடுத்தால், சுவையான புதிய பொருள்கள் எதுவும் கிடைக்காது.)

% SEGMENT OLP-0045-S08
அடுத்து, மெய்யெண்களின் மீது சில செயல்களை வரையறுக்க வேண்டும்.
ஒவ்வொரு $p \in \Rat$ க்கும் $p_\Real =
\Setabs{q \in \Rat}{q < p}$ என நிர்ணயித்து,
விகிதமுறு எண்களை மெய்யெண்களுக்குள் பதிப்பதிலிருந்து தொடங்குகிறோம்.
பின்னர் பின்வருமாறு வரையறுக்கிறோம்:
\begin{align*}
	\alpha + \beta &= \Setabs{p + q}{p \in \alpha \land q \in \beta}\\
%	\alpha - \beta &\defis \Setabs{p - q}{p \in \alpha \land q \in \Rat \setminus \beta}\\
	\alpha \times \beta &=
	\Setabs{p \times q}{0 \leq p \in \alpha \land 0 \leq q \in \beta} \cup 0^\mathbb{R} & \text{எனில் }\alpha, \beta \geq 0_\Real
%	\alpha \div \beta &\defis
%	\Setabs{\nicefrac{p}{q}}{p \in \alpha \land q \in \Rat \setminus \beta}  & \text{எனில் }\alpha \geq 0_\Real, \beta > 0_\Real
\end{align*}
பிற பெருக்கல் நிலைகளைக் கையாள, முதலில் பின்வருமாறு அமைக்க:
%$0_\Real \times \alpha = 0_\Real = \alpha \times 0_\Real$ எனவும், மேலும் பின்வருமாறும் அமைக்க:
\begin{align*}
	-\alpha &= \Setabs{p - q}{p < 0 \land q \notin \alpha}
\end{align*}
பின்னர் பின்வருமாறு நிர்ணயிக்க:
\begin{align*}
	\alpha \times \beta &\defis
	\begin{cases}
		\mathord{-}\alpha \times \mathord{-}\beta &\text{எனில் }\alpha < 0_\Real\text{ மற்றும் }\beta < 0_\Real\\
		\mathord{-}(\mathord{-}\alpha \times \beta) &\text{எனில் }\alpha < 0_\Real \text{ மற்றும் }\beta > 0_\Real\\
		\mathord{-}(\alpha \times \mathord{-}\beta) &\text{எனில் }\alpha > 0_\Real \text{ மற்றும் }\beta < 0_\Real
	\end{cases}
\end{align*}
இந்த வரையறைகள் ஒவ்வொன்றும் எப்போதும் ஒரு வெட்டைத் தருகின்றனவா
எனப் பின்னர் சரிபார்க்க வேண்டும்.
இறுதியாக, இவ்வாறு வரையறுக்கப்பட்ட வெட்டுகள்,
அவை செயல்பட வேண்டியவாறே செயல்படுகின்றன என்பதை எளிய
(ஆனால் நீண்ட) விளக்கத்தின் மூலம் சரிபார்க்க வேண்டும்.
இவை அனைத்தையும் \olref[check]{sec} இல் தருகிறோம்.
\end{document}
